% ============================================================================
% Chapter 9 solutions: Simplex Method in Tableau Form
% Book source: part1-linear-programming/ch06-simplex/simplex-tableau.tex
% Exercise numbering follows \begin{ex} order in that file (12 exercises).
% All pivots and optima verified with exact-fraction basis algebra and
% scipy.optimize.linprog, July 2026.
% ============================================================================
\section{Chapter 9: Simplex Method in Tableau Form}

Chapter 9 repackages the simplex method as row operations on a single
table. Its twelve exercises cover reading tableaus (9.1 to 9.3), executing
pivots and full solves including a Big-M setup (9.4 to 9.7), the sign
convention and the tableau-dictionary correspondence (9.8, 9.9),
recognizing unboundedness and infeasibility from a tableau (9.10, 9.11),
and one challenge problem that turns an unbounded tableau into an explicit
certificate ray (9.12). Throughout, recall the chapter's sign convention:
the $z$-row stores $z - \c^\top\x = 0$, so entering variables have
\emph{negative} $z$-row entries and the tableau is optimal when every
$z$-row entry is nonnegative. All arithmetic below was verified with
exact-fraction basis computations, and each optimum with an LP solver.
Book Selected Solutions exist for 9.2, 9.3, 9.5, and 9.6 and are
reproduced here.

\exsol{9.1}{Reading a mid-solve tableau}{1}

\textbf{(a)} The basic variables are the row labels: $s_1, s_2, y$. The
nonbasic variables are $x$ and $s_3$.

\textbf{(b)} Setting $x = s_3 = 0$ and reading the RHS column:
\[
x = 0,\quad y = 7,\quad s_1 = 2,\quad s_2 = 9,\quad s_3 = 0,
\qquad z = 21.
\]

\textbf{(c)} The $z$-row entries of the nonbasic variables are
$-\tfrac12$ (for $x$) and $\tfrac32$ (for $s_3$). In the chapter's
convention these are the recorded reduced costs; because of the stored
sign flip, increasing $x$ by one unit \emph{raises} $z$ by $\tfrac12$,
while increasing $s_3$ would lower it by $\tfrac32$.

\textbf{(d)} Not optimal: the $z$-row has the negative entry $-\tfrac12$
under $x$, so $x$ enters next. (The subsequent ratio test gives
$2/\tfrac12 = 4$, $9/\tfrac32 = 6$, $7/\tfrac12 = 14$, so $s_1$ would
leave, leading to the optimal tableau with $z = 23$.)

\exsol{9.2}{Reading a tableau}{1}

(Book Selected Solution, verified.) The basic variables are the row
labels $x_1$ and $x_2$; the nonbasic variables are $s_1$ and $s_2$.
Setting the nonbasic variables to zero and reading the RHS column gives
the basic feasible solution
\[
x_1 = 6, \quad x_2 = 4, \quad s_1 = s_2 = 0, \qquad z = 40.
\]
The objective-row entries of the nonbasic variables are $3$ and $1$, both
nonnegative. In this convention a positive entry means that increasing
the corresponding nonbasic variable would \emph{decrease} $z$, so no
pivot can improve the objective and the tableau is optimal.

\exsol{9.3}{Entering variable, ratio test, and pivot element}{1}

(Book Selected Solution, verified.) The most negative objective-row
entry is $-4$, so $x_1$ enters. The ratio test divides each RHS by the
positive entries of the $x_1$ column:
\[
\frac{6}{2} = 3, \qquad \frac{8}{4} = 2, \qquad \frac{12}{3} = 4.
\]
The minimum ratio $2$ occurs in the $s_2$ row, so $s_2$ leaves and the
pivot element is the $4$ in the $s_2$ row of the $x_1$ column.

\exsol{9.4}{Perform one pivot}{2}

Pivot on the $4$ in the $s_2$ row, $x_1$ column. Scale the pivot row by
$\tfrac14$ to get the new $x_1$-row
$[\,1,\ \tfrac14,\ 0,\ \tfrac14,\ 0 \mid 2\,]$, then clear the $x_1$
column: $s_1\text{-row} \leftarrow s_1\text{-row} - 2(x_1\text{-row})$,
$s_3\text{-row} \leftarrow s_3\text{-row} - 3(x_1\text{-row})$,
$z\text{-row} \leftarrow z\text{-row} + 4(x_1\text{-row})$. The result:
\begin{center}
\renewcommand{\arraystretch}{1.3}
\begin{tabular}{c|ccccc|c}
Basis & $x_1$ & $x_2$ & $s_1$ & $s_2$ & $s_3$ & RHS\\
\hline
$z$   & $0$ & $-2$ & $0$ & $1$ & $0$ & $8$\\
\hline
$s_1$ & $0$ & $\tfrac52$  & $1$ & $-\tfrac12$ & $0$ & $2$\\
$x_1$ & $1$ & $\tfrac14$  & $0$ & $\tfrac14$  & $0$ & $2$\\
$s_3$ & $0$ & $\tfrac{13}{4}$ & $0$ & $-\tfrac34$ & $1$ & $6$\\
\end{tabular}
\end{center}
The $x_1$ column is now the unit column $(0; 0, 1, 0)$, and the objective
value in the $z$-row RHS has increased from $0$ to $8$ (the new BFS is
$x_1 = 2$, $x_2 = 0$ with $z = 4 \cdot 2 = 8$). The tableau is not yet
optimal ($-2$ under $x_2$), but the exercise asks only for this one pivot.

\exsol{9.5}{Solve a linear program in tableau form}{2}

(Book Selected Solution, verified, with the starting tableau displayed.)
Adding slacks $s_1, s_2$, the initial tableau is
\begin{center}
\renewcommand{\arraystretch}{1.3}
\begin{tabular}{c|cccc|c|c}
Basis & $x$ & $y$ & $s_1$ & $s_2$ & RHS & ratio\\
\hline
$z$   & $-1$ & $\mathbf{-2}$ & $0$ & $0$ & $0$ & \\
\hline
$s_1$ & $1$ & $1$ & $1$ & $0$ & $4$ & $4/1 = 4$\\
$s_2$ & $1$ & $\fbox{3}$ & $0$ & $1$ & $6$ & $6/3 = 2$ (min)\\
\end{tabular}
\end{center}

\textbf{Pivot 1:} $y$ enters ($-2$ most negative); ratios $4$ and $2$,
so $s_2$ leaves; pivot on the boxed $3$. Dividing the pivot row by $3$
and clearing the $y$ column gives
\begin{center}
\renewcommand{\arraystretch}{1.3}
\begin{tabular}{c|cccc|c|c}
Basis & $x$ & $y$ & $s_1$ & $s_2$ & RHS & ratio\\
\hline
$z$   & $\mathbf{-\tfrac13}$ & $0$ & $0$ & $\tfrac23$ & $4$ & \\
\hline
$s_1$ & $\fbox{$\tfrac23$}$ & $0$ & $1$ & $-\tfrac13$ & $2$ & $2/\tfrac23 = 3$ (min)\\
$y$   & $\tfrac13$ & $1$ & $0$ & $\tfrac13$  & $2$ & $2/\tfrac13 = 6$\\
\end{tabular}
\end{center}

\textbf{Pivot 2:} $x$ enters ($-\tfrac13$); ratios $3$ and $6$, so $s_1$
leaves; pivot on the boxed $\tfrac23$:
\begin{center}
\renewcommand{\arraystretch}{1.3}
\begin{tabular}{c|cccc|c}
Basis & $x$ & $y$ & $s_1$ & $s_2$ & RHS\\
\hline
$z$   & $0$ & $0$ & $\tfrac12$ & $\tfrac12$ & $5$\\
\hline
$x$ & $1$ & $0$ & $\tfrac32$  & $-\tfrac12$ & $3$\\
$y$ & $0$ & $1$ & $-\tfrac12$ & $\tfrac12$  & $1$\\
\end{tabular}
\end{center}
Every $z$-row entry is nonnegative, so the tableau is optimal:
$(x, y) = (3, 1)$ with $z^* = 5$ (both constraints tight).

\exsol{9.6}{A full two-pivot solve}{2}

(Book Selected Solution, verified.) Adding slacks $s_1, s_2$, the
starting tableau is
\begin{center}
\renewcommand{\arraystretch}{1.3}
\begin{tabular}{c|cccc|c|c}
Basis & $x$ & $y$ & $s_1$ & $s_2$ & RHS & ratio\\
\hline
$z$   & $\mathbf{-3}$ & $-2$ & $0$ & $0$ & $0$ & \\
\hline
$s_1$ & $\fbox{2}$ & $1$ & $1$ & $0$ & $10$ & $10/2 = 5$ (min)\\
$s_2$ & $1$ & $3$ & $0$ & $1$ & $15$ & $15/1 = 15$\\
\end{tabular}
\end{center}

\textbf{Pivot 1:} $x$ enters ($-3$ most negative); $s_1$ leaves; pivot on
the boxed $2$. Scaling the pivot row by $\tfrac12$ and clearing the $x$
column:
\begin{center}
\renewcommand{\arraystretch}{1.3}
\begin{tabular}{c|cccc|c|c}
Basis & $x$ & $y$ & $s_1$ & $s_2$ & RHS & ratio\\
\hline
$z$   & $0$ & $\mathbf{-\tfrac12}$ & $\tfrac32$ & $0$ & $15$ & \\
\hline
$x$   & $1$ & $\tfrac12$ & $\tfrac12$ & $0$ & $5$ & $5/\tfrac12 = 10$\\
$s_2$ & $0$ & $\fbox{$\tfrac52$}$ & $-\tfrac12$ & $1$ & $10$ & $10/\tfrac52 = 4$ (min)\\
\end{tabular}
\end{center}

\textbf{Pivot 2:} $y$ enters ($-\tfrac12$); ratios $10$ and $4$, so
$s_2$ leaves; pivot on the boxed $\tfrac52$:
\begin{center}
\renewcommand{\arraystretch}{1.3}
\begin{tabular}{c|cccc|c}
Basis & $x$ & $y$ & $s_1$ & $s_2$ & RHS\\
\hline
$z$   & $0$ & $0$ & $\tfrac75$ & $\tfrac15$ & $17$\\
\hline
$x$ & $1$ & $0$ & $\tfrac35$ & $-\tfrac15$ & $3$\\
$y$ & $0$ & $1$ & $-\tfrac15$ & $\tfrac25$ & $4$\\
\end{tabular}
\end{center}
Every $z$-row entry is nonnegative, so the tableau is optimal:
$(x, y) = (3, 4)$ with $z^* = 17$ and $s_1 = s_2 = 0$ (both constraints
are tight at the optimum).

\exsol{9.7}{Setting up a Big-M tableau}{2}

The augmented constraints are
\[
x + y - e_1 + a_1 = 4,
\qquad
x + 2y + s_2 = 10,
\qquad
x, y, e_1, s_2, a_1 \geq 0,
\]
with objective $z - 4x - 3y + M a_1 = 0$. Written naively, the tableau is
\begin{center}
\renewcommand{\arraystretch}{1.3}
\begin{tabular}{c|ccccc|c}
Basis & $x$ & $y$ & $e_1$ & $s_2$ & $a_1$ & RHS\\
\hline
$z$   & $-4$ & $-3$ & $0$ & $0$ & $M$ & $0$\\
\hline
$a_1$ & $1$ & $1$ & $-1$ & $0$ & $1$ & $4$\\
$s_2$ & $1$ & $2$ & $0$  & $1$ & $0$ & $10$\\
\end{tabular}
\end{center}
This is not yet a proper tableau: the \emph{basic} variable $a_1$ has a
nonzero entry ($M$) in the $z$-row. Pricing out
($z\text{-row} \leftarrow z\text{-row} - M \cdot a_1\text{-row}$) gives
\begin{center}
\renewcommand{\arraystretch}{1.3}
\begin{tabular}{c|ccccc|c|c}
Basis & $x$ & $y$ & $e_1$ & $s_2$ & $a_1$ & RHS & ratio\\
\hline
$z$   & $\mathbf{-4-M}$ & $-3-M$ & $M$ & $0$ & $0$ & $-4M$ & \\
\hline
$a_1$ & $\fbox{1}$ & $1$ & $-1$ & $0$ & $1$ & $4$ & $4/1 = 4$ (min)\\
$s_2$ & $1$ & $2$ & $0$  & $1$ & $0$ & $10$ & $10/1 = 10$\\
\end{tabular}
\end{center}
Now $a_1$'s objective entry is zero, as required. The most negative
$z$-row entry is $-4 - M$ (for large $M$ it beats $-3 - M$), so
\emph{$x$ enters first}. The ratio test gives $4/1 = 4$ in the $a_1$ row
versus $10/1 = 10$ in the $s_2$ row, so \emph{yes: the artificial
variable $a_1$ leaves on the very first pivot}, and the method continues
with an ordinary feasible tableau (the pivot lands at $(x, y) = (4, 0)$
with $z = 16$).

\exsol{9.8}{The $z$-row sign convention}{2}

\textbf{(a)} A dictionary displays $z = \bar z + \sum_j \bar c_j x_j$ over
the nonbasic variables. The tableau instead stores this equation with
everything on the left: $z - \sum_j \bar c_j x_j = \bar z$. Each
coefficient $\bar c_j$ crosses the equals sign and gets negated, so a
variable that is attractive in the dictionary (positive $\bar c_j$) shows
up in the tableau with a negative entry $-\bar c_j$. The entering rules
``most positive dictionary coefficient'' and ``most negative $z$-row
entry'' therefore select the same variable.

\textbf{(b)} From $z = 12 + 3x_1 - 2x_2$, move the variable terms left:
$z - 3x_1 + 2x_2 = 12$. The $z$-row is
\[
\begin{array}{c|cc|c}
 & x_1 & x_2 & \text{RHS} \\
\hline
z & -3 & 2 & 12
\end{array}
\]
(entries $0$ under any slack or basic columns not shown).

\textbf{(c)} The stored equation is $z - \sum_j \bar c_j x_j = \bar z$,
and it only involves \emph{nonbasic} variables besides $z$. At the
current basic solution every nonbasic variable equals zero, so the
equation collapses to $z = \bar z$: the RHS entry is the current
objective value. The constant term never crosses the equals sign (it was
already on the right), so it is the one number in the row that is not
negated.

\exsol{9.9}{From tableau to dictionary}{2}

Each tableau row is an equation. The $x_1$ row says
$x_1 + 2s_1 - s_2 = 6$ and the $x_2$ row says $x_2 - s_1 + s_2 = 4$;
the $z$-row says $z + 3s_1 + s_2 = 40$. Solving each for its basic
variable (and $z$):
\[
\begin{aligned}
\max\quad & z = 40 - 3s_1 - s_2 \\
\st\quad & x_1 = 6 - 2s_1 + s_2, \\
& x_2 = 4 + s_1 - s_2, \\
& x_1, x_2, s_1, s_2 \geq 0.
\end{aligned}
\]
Setting the nonbasic variables $s_1 = s_2 = 0$ gives $x_1 = 6$,
$x_2 = 4$, $z = 40$: the same basic feasible solution as the tableau.
(Both dictionary coefficients in the objective are negative, matching the
optimality of the tableau.)

\exsol{9.10}{Recognizing an unbounded problem}{2}

The only negative $z$-row entry is $-2$ under $x_1$, so $x_1$ should
enter. The ratio test looks for \emph{positive} entries in the $x_1$
column, but they are $-1$ (in the $x_2$ row) and $-3$ (in the $s_1$
row): there are none, so the ratio test fails; no row limits the
increase of $x_1$. Increasing $x_1 = t$ while keeping $s_2 = 0$ gives
$x_2 = 3 + t$ and $s_1 = 5 + 3t$, which stay nonnegative for all
$t \geq 0$, while $z = 10 + 2t \to \infty$. The linear program is
\textbf{unbounded}: it has feasible solutions of arbitrarily large
objective value and no optimal solution.

\exsol{9.11}{Reading off infeasibility}{2}

Yes, the tableau is optimal for the penalized (Big-M) problem: for large
$M$ every $z$-row entry ($0, 0, M, 2, 0$) is nonnegative, so no pivot can
improve the penalized objective. But the artificial variable is still
basic at the positive value $a_1 = 3$ (read from its RHS). A feasible
solution of the original problem would allow all artificial variables to
be set to zero, so if the best the penalized problem can do still keeps
$a_1 = 3 > 0$, no feasible solution of the original program exists: the
original linear program is \textbf{infeasible}.

The objective value $8 - 3M$ tells the same story: it contains the
unavoidable penalty $-M a_1 = -3M$. Since $M$ stands for an arbitrarily
large constant, $8 - 3M$ is arbitrarily bad; the penalized problem cannot
escape the penalty, which is exactly the algebraic signature of
infeasibility.

\exsol{9.12}{Certifying unboundedness from a tableau}{3}

\textbf{(a)} The only negative $z$-row entry is $-2$ under $x_2$, so
$x_2$ enters. Its column entries are $-1$ (in the $x_1$ row) and $0$
(in the $s_2$ row); the ratio test requires a \emph{positive} entry and
finds none, so no row limits the increase of $x_2$: the problem is
unbounded.

\textbf{(b)} Set $x_2 = t \geq 0$ and keep the other nonbasic variable
$s_1 = 0$. The $x_1$ row ($x_1 - x_2 + s_1 = 1$) gives $x_1 = 1 + t$,
and the $s_2$ row ($s_1 + s_2 = 3$) gives $s_2 = 3$. So
\[
\x(t) = (x_1, x_2, s_1, s_2) = (1 + t,\ t,\ 0,\ 3), \qquad t \geq 0.
\]
Direct substitution into the original constraints:
\[
x_1 - x_2 = (1 + t) - t = 1 \leq 1 \quad\checkmark
\qquad
-x_1 + x_2 = -(1 + t) + t = -1 \leq 2 \quad\checkmark
\]
and $x_1 = 1 + t \geq 0$, $x_2 = t \geq 0$, $s_1 = 0 \geq 0$,
$s_2 = 2 - (-1) = 3 \geq 0$. So $\x(t)$ is feasible for every
$t \geq 0$.

\textbf{(c)} Along the ray, $z(t) = x_1 + x_2 = 1 + 2t \to \infty$ as
$t \to \infty$. Geometrically, the feasible region is the strip between
the parallel-ish lines $x_1 - x_2 = 1$ and $-x_1 + x_2 = 2$ in the first
quadrant; it is unbounded in the direction $\mathbf{d} = (1, 1)$. The
certificate ray starts at the vertex $(1, 0)$ and points along
$(1, 1)$, hugging the boundary $x_1 - x_2 = 1$; since
$\c^\top \mathbf{d} = 1 + 1 = 2 > 0$, the objective grows without bound
along it.
